\begin{lemma}[The gradient sup bound in the nondegenerate case]
Let $n\ge2$, $p>1$ and $r>0$, suppose the Sobolev--Poincar\'e inequality on
balls holds in dimension $n$ (\noderef{CitedSobolevPoincareBall}), and assume
the cited regularity result (C1) (\noderef{CitedRegularityEpsPos}) for
\emph{every} field $A'\colon\R^n\to\R^n$, \emph{every} open set
$\Omega'\subseteq\R^n$ and all parameters $\varepsilon',\gamma_0',\gamma_1'$
with $\gamma_0'>0$ and $\gamma_1'>0$.  Then there is a constant $C=C(p,r,n)>0$
with the following property.  Let $\varepsilon>0$, $\gamma_0>0$ and $\gamma_1>0$
with $\gamma_1\le r\gamma_0$, let $\Omega\subseteq\R^n$ be open, let
$A\colon\R^n\to\R^n$ satisfy (Q1) of the thesis--(Q2) of the thesis on $\Omega$ with
parameters $p,\varepsilon,\gamma_0,\gamma_1$ (\noderef{EllipticStructureQ12}),
let $u$ be a weak solution of $\operatorname{div}A(\nabla u)=0$ in $\Omega$
(\noderef{IsWeakSolutionOn}) with weak gradient $g$, and let $y\in\R^n$ and
$\rho>0$ satisfy $\overline B_\rho(y)\subseteq\Omega$.  Then for almost every
$x\in B_{\rho/2}(y)$,
\[
\bigl(\varepsilon+|g(x)|^2\bigr)^{\frac p2}
\ \le\ \frac{C}{\rho^n}\int_{B_\rho(y)}\bigl(\varepsilon+|g|^2\bigr)^{\frac p2}dx .
\]
The conclusion is an almost-everywhere statement about the given weak gradient
$g$, and deliberately so: a weak gradient is an almost-everywhere class, so no
pointwise assertion may be made about $g$ itself.  The constant $C$ is bound
before $\varepsilon,\gamma_0,\gamma_1,\Omega,A,u,g,y,\rho$, and the structure
constants enter only through the bound $\gamma_1\le r\gamma_0$; no restriction
$\rho=1$ is imposed and no rescaling is needed, the two quantitative inputs
being already scale-invariant.
\end{lemma}
\begin{proof}
Write $B_\sigma(y)$ for the open ball of radius $\sigma$ about $y$,
$\overline B_\sigma(y)$ for the closed one, $|E|$ for Lebesgue measure, $D_i$
for the derivative in the direction of the $i$-th standard basis vector, and
$(\cdot,\cdot)$, $|\cdot|$ for the Euclidean inner product and norm; all
non-integer powers are real powers of nonnegative numbers.

\emph{Step 0: the constant, chosen before all the data.}  Let $c>0$ be the
constant that \noderef{CaccioppoliPowerWeight} supplies for the triple
$(n,p,r)$; its hypotheses there are $2\le n$, $1<p$ and $r>0$, all of which are
hypotheses of the present lemma, so $c=c(n,p,r)$.  Let $C_0>0$ be the constant
that \noderef{MoserLocalBoundedness} supplies for the pair $(n,Q)$ with $Q=c$;
its hypotheses there are $2\le n$, $0<Q$ and the Sobolev--Poincar\'e inequality
on balls in dimension $n$ (\noderef{CitedSobolevPoincareBall}), which is a
hypothesis of the present lemma.  So $C_0=C_0(n,c)$.  Put
\[
C=C_0^{\,2}>0 .
\]
Both $c$ and $C_0$, hence $C$, are determined by $n,p,r$ alone and are fixed
before $\varepsilon,\gamma_0,\gamma_1,\Omega,A,u,g,y,\rho$, which is the
quantifier order the statement requires.

Now fix data as in the statement.

\emph{A1: what (C1) supplies.}  The present lemma assumes the cited regularity
result (C1) for every field, every open set and all parameters with positive
structure constants.  Instantiate that clause at
$\varepsilon,\gamma_0,\gamma_1,\Omega,A$: its side conditions $\gamma_0>0$,
$\gamma_1>0$ and openness of $\Omega$ hold.  The resulting predicate
\noderef{CitedRegularityEpsPos} has as its own antecedents $\varepsilon>0$, the
structure conditions on $\Omega$ (\noderef{EllipticStructureQ12}) and the
weak-solution property (\noderef{IsWeakSolutionOn}), all of which hold, and it
yields
\begin{enumerate}
\item $u\in W^{1,2}_{\mathrm{loc}}(\Omega)$ with weak gradient $g$
  (\noderef{MemW1pLocOn}); and
\item a field $\hat g$ with (i) $\hat g$ continuous on $\Omega$, (ii)
  $g=\hat g$ almost everywhere on $\Omega$, and (iii) a family
  $D\hat g=(D\hat g^1,\dots,D\hat g^n)$ with
  $\hat g^k\in W^{1,2}_{\mathrm{loc}}(\Omega)$ and weak gradient $D\hat g^k$ for
  each $k$ (\noderef{MemW1pLocOn}).
\end{enumerate}
Fix such a $\hat g$ and such a family $D\hat g$ once and for all.

\emph{A warning about the two fields.}  In this proof $g$ is the lemma's own
given field --- the weak gradient of $u$, named in the hypotheses --- and
$\hat g$ is the representative that (C1) produces.  Clauses (i) and (iii) are
properties of $\hat g$ and \emph{not} of $g$: a weak gradient is an
almost-everywhere class, so neither continuity nor pointwise boundedness can be
asserted of $g$, and \noderef{CitedRegularityEpsPos} is stated with the
representative existentially bound for exactly that reason.  \emph{The
consequence is structural, not cosmetic: $w$ and $v$ below are built from
$\hat g$, steps A2--A6 are run for $\hat g$, and the passage back to $g$ is made
once, at A7, where the conclusion is already an almost-everywhere statement.}
This is available because $w$ and $v$ are internal to this proof --- no step and
no cited node fixes their names --- so the representative may be carried where
points are needed and the given field where only integrals are.

\emph{A2: the functions $w$ and $v$, built from the representative.}  Put
\[
w=\varepsilon+|\hat g|^2,\qquad v=w^{\frac p4},
\]
both defined at every point.  On $\Omega$ we have $w\ge\varepsilon>0$, so $v>0$
there; and $w$ is continuous on $\Omega$ by A1(2)(i), hence so is $v$, because
$s\mapsto s^{p/4}$ is continuous on $(0,\infty)$.  Note for later that
$v^2=w^{p/2}$ and, for $\theta>0$, $(v^\theta)^2=w^{p\theta/2}$, the exponent
rules for real powers being available since $w>0$ on $\Omega$.

\emph{A3: $v^\theta$ is a local $W^{1,2}$ function, with an explicit weak
gradient.}  Fix $t>0$ and put $\theta=1+\frac tp$, so that $\theta>1$ and
$v^\theta=w^{\frac{p+t}4}$; as $t$ runs over $(0,\infty)$, $\theta$ runs over
all of $(1,\infty)$, and $t=p(\theta-1)$ recovers $t$ from $\theta$.  Apply
\noderef{WeakGradientVectorComp} with $n$, $m=n$, $q=2$, the open set $\Omega$,
the functions $G_k=\hat g^k$ with weak gradients $DG_k=D\hat g^k$, and
\[
F(y_1,\dots,y_n)=\Bigl(\varepsilon+\sum_{k=1}^ny_k^2\Bigr)^{\frac{p+t}4}.
\]
Its hypotheses hold: $1\le n$ and $1\le2$; $\Omega$ is open; each
$G_k\in W^{1,2}_{\mathrm{loc}}(\Omega)$ with weak gradient $DG_k$ by A1(2)(iii);
the family is locally uniformly bounded, since for compact $K\subseteq\Omega$
the function $|\hat g|$ is continuous on $K$ by A1(2)(i) and therefore bounded
there, say by $L$, whence $|G_k(x)|\le|\hat g(x)|\le L$ for all $x\in K$ and all
$k$.  \emph{This is the hypothesis that forces the representative on us: it asks
for a bound at \emph{every} point of $K$, and the given $g$ need not satisfy
one, a weak gradient being alterable on a null set without disturbing any
hypothesis of this lemma.}  And $F$ is of class $C^1$ on $\R^n$, being the
composition of the polynomial $y\mapsto\varepsilon+|y|^2$, which is smooth and
takes values in $[\varepsilon,\infty)\subseteq(0,\infty)$ because
$\varepsilon>0$, with $s\mapsto s^{(p+t)/4}$, which is smooth on $(0,\infty)$.
Since $F(G(x))=\bigl(\varepsilon+|\hat g(x)|^2\bigr)^{(p+t)/4}=v^\theta(x)$, the
conclusion is that $v^\theta\in W^{1,2}_{\mathrm{loc}}(\Omega)$ with weak
gradient
\[
G_\theta=\sum_{k=1}^nD_kF(G)\,D\hat g^k
=\frac{p+t}4\,w^{\frac{p+t}4-1}\,\nabla w,
\qquad
\nabla w:=2\sum_{k=1}^n\hat g^k\,D\hat g^k ,
\]
where the middle equality is the computation
$D_kF(y)=\frac{p+t}4\bigl(\varepsilon+|y|^2\bigr)^{\frac{p+t}4-1}\cdot2y_k$
evaluated at $y=G(x)$.  Finally $B_\rho(y)$ is open and contained in
$\overline B_\rho(y)\subseteq\Omega$, so clause (2) of
\noderef{WeakGradientRestrict}, applied with $q=2$, the sets
$\Omega\supseteq B_\rho(y)$ and the pair $(v^\theta,G_\theta)$, gives
\[
v^\theta\in W^{1,2}_{\mathrm{loc}}\bigl(B_\rho(y)\bigr)
\quad\text{with weak gradient}\quad G_\theta .
\tag{A3}
\]

\emph{A4: the Caccioppoli inequality in the form Moser's lemma consumes.}
Apply \noderef{CaccioppoliPowerWeight} for the triple $(n,p,r)$, whose constant
$c$ was fixed in Step 0, to the data
$\varepsilon,\gamma_0,\gamma_1,t,\Omega,A,u,g,D\hat g$ and an arbitrary $\zeta$
which is of class $C^\infty$ on $\R^n$, has compact support, and satisfies
$\mathrm{tsupport}\,\zeta\subseteq\Omega$ --- these are the three conditions
that node binds on its cut-off, and $\mathrm{tsupport}$, the \emph{closure} of
the support, is the one it asks about, so it is the one used here and below.
Its hypotheses are $\varepsilon>0$, $\gamma_0>0$, $\gamma_1\le r\gamma_0$,
$t>0$, openness of $\Omega$, the structure conditions
(\noderef{EllipticStructureQ12}), the cited regularity result (C1)
(\noderef{CitedRegularityEpsPos}) at these very data, which is the instance A1
already formed, the weak-solution property (\noderef{IsWeakSolutionOn}), the
membership $g^k\in W^{1,2}_{\mathrm{loc}}(\Omega)$ with weak gradient
$D\hat g^k$ for every $k$ (\noderef{MemW1pLocOn}), established in the next
paragraph, and the three cut-off conditions just listed.

\emph{The one transfer back to the given field.}  That node carries its
gradient-regularity hypothesis on the \emph{same} field as its weak-solution
hypothesis, namely on $g$, whereas A1(2)(iii) supplies it for $\hat g$.  The two
are reconciled once, and the fact needed is that $\mathrm{MemW1pLocOn}$ is
unchanged when the function is modified on a null set, the weak gradient being
kept.  By \noderef{MemW1pLocOn} that assertion is two things, and each is an
integral one.  First the identity of \noderef{HasWeakGradientOn}: for every test
function $\psi$ of that node's class and every index $i$,
$\int_\Omega\hat g^k\,D_i\psi\,dx=-\int_\Omega(D\hat g^k)_i\,\psi\,dx$, and
$g^k=\hat g^k$ almost everywhere on $\Omega$ by A1(2)(ii), so the two left-hand
integrands agree almost everywhere and the integrals are equal, giving
$\int_\Omega g^k\,D_i\psi\,dx=-\int_\Omega(D\hat g^k)_i\,\psi\,dx$.  Second the
$L^2$ clauses: for compact $K\subseteq\Omega$ we have $g^k=\hat g^k$ almost
everywhere on $K$ and $D\hat g^k$ is not changed at all, so $g^k\in L^2(K)$ and
$D\hat g^k\in L^2(K)$ follow from the same statements for $\hat g^k$.  Hence
\[
g^k\in W^{1,2}_{\mathrm{loc}}(\Omega)\ \text{ with weak gradient }\ D\hat g^k
\qquad(k=1,\dots,n),
\tag{A4a}
\]
which is the hypothesis just invoked.  (No further node is needed for it: the
two clauses of \noderef{MemW1pLocOn} are the only facts used, and both are
integral assertions.)

Its conclusion is
\[
\int_\Omega\Bigl(\frac{p+t}4\Bigr)^2w_\flat^{\frac{p+t}2-2}
\Bigl(\sum_{i=1}^n(\nabla w_\flat)_i^2\Bigr)\zeta^2\,dx
\ \le\ c\int_\Omega w_\flat^{\frac{p+t}2}\,|\nabla\zeta|^2\,dx ,
\]
where $w_\flat=\varepsilon+|g|^2$ and
$(\nabla w_\flat)_i=2\sum_kg^k(D\hat g^k)^i$ are the weight and the field that
node forms from the data it was given, i.e.\ from $g$ and $D\hat g$.

\emph{The weights agree almost everywhere, so the display may be rewritten in
$w$.}  By A1(2)(ii), $g=\hat g$ almost everywhere on $\Omega$, and $w_\flat$,
$\nabla w_\flat$ are obtained from $g$ by the same two formulas that produce
$w$, $\nabla w$ from $\hat g$ --- the second with the \emph{same} family
$D\hat g$ in both cases --- so
\[
w_\flat=w\quad\text{and}\quad\nabla w_\flat=\nabla w
\qquad\text{almost everywhere on }\Omega ,
\tag{A4b}
\]
and hence so does every real power of the weight.  Both sides of the display are
integrals, so altering their integrands on a null set changes neither, and the
display holds verbatim with $w$ and $\nabla w$ in place of $w_\flat$ and
$\nabla w_\flat$.  Now by A3
\[
|G_\theta|^2=\Bigl(\frac{p+t}4\Bigr)^2\bigl(w^{\frac{p+t}4-1}\bigr)^2|\nabla w|^2
=\Bigl(\frac{p+t}4\Bigr)^2w^{\frac{p+t}2-2}\sum_{i=1}^n(\nabla w)_i^2
\]
on $\Omega$, using $(w^a)^2=w^{2a}$ for $w>0$ and
$2\bigl(\frac{p+t}4-1\bigr)=\frac{p+t}2-2$; and
$(v^\theta)^2=w^{p\theta/2}=w^{\frac{p+t}2}$ by A2.  So the display reads
\[
\int_\Omega|G_\theta|^2\zeta^2\,dx
\ \le\ c\int_\Omega(v^\theta)^2|\nabla\zeta|^2\,dx
\]
for every such $\zeta$.  If moreover $\mathrm{tsupport}\,\zeta\subseteq
B_\rho(y)$, then both integrands vanish at every point outside
$\mathrm{tsupport}\,\zeta$: there $\zeta=0$, because
$\mathrm{supp}\,\zeta\subseteq\mathrm{tsupport}\,\zeta$, so $\zeta^2=0$; and the
complement of $\mathrm{tsupport}\,\zeta$ is open with $\zeta$ identically $0$ on
it, so $\nabla\zeta=0$ there as well.  Hence each integrand has the same product
with the indicator of $\Omega$ as with the indicator of $B_\rho(y)$, the two
indicators agreeing on $\mathrm{tsupport}\,\zeta$, and the two integrals may be
taken over $B_\rho(y)$:
\[
\int_{B_\rho(y)}|G_\theta|^2\zeta^2\,dx
\ \le\ c\int_{B_\rho(y)}(v^\theta)^2|\nabla\zeta|^2\,dx
\tag{A4}
\]
for every $\zeta$ of class $C^\infty$ with compact support and
$\mathrm{tsupport}\,\zeta\subseteq B_\rho(y)$ --- which is exactly the cut-off
class the Caccioppoli hypothesis of \noderef{MoserLocalBoundedness} quantifies
over, at the ball $B_\rho(y)$.  The constant $c$ does not depend on $t$, hence
not on $\theta$: it was chosen in Step 0, before $t$ was introduced, and this
uniformity is exactly what \noderef{MoserLocalBoundedness} requires of its $Q$.

\emph{A5: truncation, and the three remaining hypotheses of Moser's lemma.}  Put
$\tilde v=\mathbf1_{B_\rho(y)}v$, the function equal to $v$ on $B_\rho(y)$ and
to $0$ elsewhere.  Then:
\begin{enumerate}
\item $\tilde v\ge0$ everywhere, since $v=w^{p/4}\ge0$ as a real power of a
  nonnegative number.
\item $\tilde v$ is measurable.  For $a<0$ the superlevel set $\{\tilde v>a\}$
  is all of $\R^n$; for $a\ge0$ it is $\{x\in B_\rho(y):v(x)>a\}$, which is open
  because $v$ is continuous on the open set $B_\rho(y)\subseteq\Omega$ by A2 ---
  legitimately, $v$ having been built from the continuous $\hat g$ and not from
  $g$.  All superlevel sets being measurable, $\tilde v$ is measurable
  (Preamble).
\item $\tilde v^2$ is integrable on $B_\rho(y)$.  Indeed $\tilde v^2=w^{p/2}$
  there, $\overline B_\rho(y)$ is a compact subset of $\Omega$ and $w$ is
  continuous on $\Omega$, so $w^{p/2}$ is continuous on $\overline B_\rho(y)$
  and therefore bounded there, while $|B_\rho(y)|<\infty$; the Preamble item on
  domination gives integrability.
\end{enumerate}
Moreover $\tilde v^\theta=v^\theta$ on $B_\rho(y)$ for every $\theta>0$.  The
two hypotheses of \noderef{MoserLocalBoundedness} that mention the powers of the
function are statements about $B_\rho(y)$ only: the defining identity of
\noderef{HasWeakGradientOn} on $B_\rho(y)$ and the local $L^2$ clauses of
\noderef{MemW1pLocOn} on $B_\rho(y)$ integrate over $B_\rho(y)$ and over its
compact subsets, and (A4) integrates over $B_\rho(y)$.  Since $\tilde v^\theta$
and $v^\theta$ agree there, (A3) and (A4) hold verbatim with $\tilde v$ in place
of $v$.

\emph{A6: Moser's lemma.}  Apply \noderef{MoserLocalBoundedness} with $n$,
$Q=c$, the hypothesis \noderef{CitedSobolevPoincareBall} of the present lemma,
and then with the centre $y$, the radius $\rho$, the function $\tilde v$ and the
family of fields $\theta\mapsto G_{p(\theta-1)}$ for $\theta>1$ (and, say, $0$
for $\theta\le1$; the hypotheses only ever evaluate it at $\theta>1$).  Its
hypotheses are: $2\le n$ and $0<c$ (Step 0); (C6), which the present lemma
assumes; $0<\rho$; measurability and nonnegativity of $\tilde v$ and
integrability of $\tilde v^2$ on $B_\rho(y)$, which are A5(1)--(3); the
membership $\tilde v^\theta\in W^{1,2}_{\mathrm{loc}}(B_\rho(y))$ with weak
gradient $G_{p(\theta-1)}$ for every $\theta>1$, which is (A3) together with A5,
since $t=p(\theta-1)>0$ and then $\theta=1+t/p$; and the Caccioppoli inequality
for every $\theta>1$ and every cut-off, with the single constant $c$, which is
(A4) together with A5.  The conclusion is: for almost every
$x\in B_{\rho/2}(y)$,
\[
\tilde v(x)\ \le\ C_0\Bigl(\rho^{-n}\int_{B_\rho(y)}\tilde v^2\,dx\Bigr)^{\frac12}.
\]

\emph{A7: squaring, and the passage back to the given field.}  Let
$x\in B_{\rho/2}(y)$ be a point where the last display holds \emph{and} at which
$g(x)=\hat g(x)$; almost every point of $B_{\rho/2}(y)$ is such a point, being
in the intersection of two sets of full measure in $B_\rho(y)$ --- the first
from A6's conclusion, the second from A1(2)(ii).  There
$\tilde v(x)=v(x)=w(x)^{p/4}$, since $B_{\rho/2}(y)\subseteq B_\rho(y)$, and
$w(x)=\varepsilon+|\hat g(x)|^2=\varepsilon+|g(x)|^2$ by the choice of $x$.
Both sides of the display are nonnegative, so squaring preserves the inequality,
and $\bigl(w(x)^{p/4}\bigr)^2=w(x)^{p/2}$ while
$\int_{B_\rho(y)}\tilde v^2=\int_{B_\rho(y)}w^{p/2}$ by A5(3); and
$\int_{B_\rho(y)}w^{p/2}=\int_{B_\rho(y)}(\varepsilon+|g|^2)^{p/2}$, the two
integrands agreeing almost everywhere by (A4b).  \emph{This is the single point
at which the argument returns to $g$, and it is legitimate because the
conclusion is itself an almost-everywhere statement and its right-hand side is
an integral.}  Hence for almost every $x\in B_{\rho/2}(y)$
\[
\bigl(\varepsilon+|g(x)|^2\bigr)^{\frac p2}=w(x)^{\frac p2}
\ \le\ \frac{C_0^{\,2}}{\rho^n}\int_{B_\rho(y)}w^{\frac p2}\,dx
=\frac{C}{\rho^n}\int_{B_\rho(y)}\bigl(\varepsilon+|g|^2\bigr)^{\frac p2}dx ,
\]
which is the assertion, with the constant $C=C_0^{\,2}$ of Step 0.
\end{proof}
